\rev{Physics-informed neural networks that output the spectral coefficients of the solution compute spatial derivatives exactly. We evaluate such a Chebyshev-coefficient network for the one-dimensional Allen--Cahn equation with Neumann boundary conditions on three benchmarks, with five seeds, a certified reference solution and a classical Chebyshev solver as baseline. An exactly Neumann-compatible basis keeps the network accurate at large mode numbers, where a soft boundary penalty fails. Trained with Adam, the network is \N{b1.ratio.van_fhK32} and \N{b2.ratio.van_hard} times more accurate than a space--time PINN on a reaction-dominated and a diffusion-active benchmark; after L-BFGS refinement the advantage shrinks or vanishes. The classical solver is more accurate still and far cheaper, and on a long-horizon benchmark with moving interfaces both networks fail. Every trained coefficient network has monotone energy, so an energy penalty is inactive, while the maximum bound principle is not guaranteed. An a-posteriori energy balance, machine-checked in Lean~4, bounds the energy growth of any trained trajectory by its measured residual. In the inverse problem, joint training returns an interface parameter set by its starting value; variable projection removes this on noise-free data but is biased under noise, where classical least squares is better. A reflection-symmetric Fourier basis preserves two-dimensional symmetry to round-off.}
